\title{Chain-of-Thought Prompting Elicits Reasoning in Large Language Models}

\begin{abstract}
We investigate how producing a \textit{chain of thought}—a sequence of intermediate reasoning steps—substantially enhances the capability of large language models to handle complex reasoning tasks. Specifically, we demonstrate that such reasoning skills emerge naturally in sufficiently large language models through a straightforward approach termed \textit{chain-of-thought prompting}, which involves providing a few chain of thought examples as prompts.

Our experiments involving three large language models reveal that chain-of-thought prompting boosts performance across various arithmetic, commonsense, and symbolic reasoning challenges. The improvements observed can be remarkable. For example, prompting a PaLM 540B with only eight chain-of-thought exemplars attains state-of-the-art results on the GSM8K benchmark for math word problems, outperforming even finetuned GPT-3 equipped with a verifier.
\end{abstract}

\section{Introduction}
The field of NLP has recently undergone a transformation due to the advent of language models \citep[][\textit{inter alia}]{peters-etal-2018-deep,devlin-etal-2019-bert,brown2020language}. Increasing the size of language models has been shown to offer numerous advantages, including better performance and improved sample efficiency \citep[\textit{inter alia}]{kaplan2020scaling,brown2020language}. Nonetheless, merely scaling up model size has not been enough to achieve strong results on demanding tasks such as arithmetic, commonsense, and symbolic reasoning \citep{rae2021scaling}.

This paper examines how a large language model’s reasoning capacity can be unlocked using a simple technique inspired by two principles. First, methods for arithmetic reasoning benefit from generating natural language explanations that lead to the final solution. Previous research has enabled models to produce natural language intermediate steps either by training from scratch \citep{ling-etal-2017-program} or by finetuning pretrained models \citep{cobbe2021training}, as well as through neuro-symbolic methods that employ formal languages instead of natural language \citep{roy-roth-2015-solving, chiang-chen-2019-semantically,amini-etal-2019-mathqa, chen2019neural}. Second, large language models present an appealing opportunity for in-context few-shot learning via \emph{prompting}. Rather than finetuning a separate model checkpoint for each new task, one can simply provide the model with a few input-output examples that illustrate the task. Impressively, this approach has been successful for various straightforward question-answering tasks \citep{brown2020language}.

Nevertheless, both of the aforementioned approaches have significant drawbacks. For rationale-augmented training and fine-tuning techniques, generating a substantial collection of high-quality rationales is expensive and far more complex than producing simple input-output pairs typically used in standard machine learning. Meanwhile, the conventional few-shot prompting method employed in \citet{brown2020language} tends to perform poorly on tasks that demand reasoning skills and often shows limited improvement even as the language model size increases \citep{rae2021scaling}. In this study, we merge the advantages of these two strategies while circumventing their respective limitations. Specifically, we investigate the capability of language models to carry out few-shot prompting for reasoning tasks by using prompts composed of triples: $\langle$input, \emph{chain of thought}, output$\rangle$. A \emph{chain of thought} consists of a sequence of intermediate natural language reasoning steps that culminate in the final output; we term this method \textit{chain-of-thought prompting}. An illustrative example of such a prompt is provided in \cref{fig:pull-figure}.

We conduct empirical tests on benchmarks involving arithmetic, commonsense, and symbolic reasoning, demonstrating that chain-of-thought prompting surpasses standard prompting, occasionally by a substantial margin. One such outcome is depicted in \cref{fig:pull-bar-chart}—on the GSM8K benchmark for math word problems \citep{cobbe2021training}, chain-of-thought prompting using \palm{} 540B significantly outperforms standard prompting and sets a new state-of-the-art record. Employing a prompting-only approach is crucial because it obviates the need for a large training dataset and enables a single model checkpoint to handle multiple tasks without losing generality. This research highlights how large language models can effectively learn from a few examples that include natural language information about the task (as opposed to learning patterns from extensive training datasets linking inputs to outputs).

\section{Chain-of-Thought Prompting} When tackling a complex reasoning task like a multi-step math word problem, it is common to break down the problem into intermediate steps and address each one before arriving at the final solution: \textit{``After Jane gives 2 flowers to her mom she has 10 $\ldots$ then after she gives 3 to her dad she will have 7 $\ldots$ so the answer is 7.''} This paper aims to equip language models with the capability to produce a similar \textit{chain of thought}—a logical sequence of intermediate reasoning steps culminating in the final answer to the problem. We demonstrate that sufficiently large language models can generate such chains of thought when provided with chain-of-thought reasoning examples in few-shot prompting.

\cref{fig:pull-figure} presents an instance where a model generates a chain of thought to solve a math word problem it would have otherwise answered incorrectly. The chain of thought here resembles a solution and can be interpreted as such; however, we prefer to refer to it as a chain of thought to emphasize that it emulates a stepwise cognitive process leading to the answer (and also because solutions or explanations typically follow the final answer \citep[][\textit{inter alia}]{narang2020wt5,wiegreffe2021reframing,lampinen2022can}).

Chain-of-thought prompting offers several appealing characteristics as a method to enhance reasoning capabilities in language models.

\begin{enumerate}[topsep=1pt,itemsep=0ex]
    \item Firstly, chain of thought inherently enables models to break down complex multi-step problems into simpler intermediate steps, allowing additional computational resources to be focused on problems requiring extended reasoning.
    \item Secondly, a chain of thought offers a transparent view of the model’s decision-making process, illustrating how a specific conclusion might have been reached and providing a means to identify errors along the reasoning path (although fully understanding the model’s internal computations leading to an answer remains an unresolved issue).
    \item Thirdly, chain-of-thought reasoning can be applied to tasks such as mathematical word problems, commonsense inference, and symbolic manipulation, and is potentially useful (at least in theory) for any problem that humans solve through language.
    \item Lastly, chain-of-thought reasoning can be easily triggered in sufficiently large pre-trained language models by simply including chain-of-thought examples within the few-shot prompt exemplars.
\end{enumerate}

In our empirical evaluations, we demonstrate the effectiveness of chain-of-thought prompting across arithmetic reasoning (\cref{sec:arithmetic-reasoning}), commonsense reasoning (\cref{sec:commonsense-reasoning}), and symbolic reasoning tasks (\cref{sec:symbolic-reasoning}).

\section{Arithmetic Reasoning}\label{sec:arithmetic-reasoning}
We start by examining math word problems as shown in \cref{fig:pull-figure}, which assess language models’ arithmetic reasoning skills. Although these problems are straightforward for humans, language models frequently find arithmetic reasoning challenging \citep[][\textit{inter alia}]{hendrycks2021measuring,patel-etal-2021-nlp}. Remarkably, when chain-of-thought prompting is employed with the 540B parameter model, performance matches that of task-specific fine-tuned models on several datasets, even setting new state-of-the-art results on the difficult GSM8K benchmark \citep{cobbe2021training}.

\subsection{Experimental Setup}

We investigate chain-of-thought prompting across different language models on a variety of benchmarks.

\textbf{Benchmarks.} We utilize five math word problem datasets:
\textbf{(1)} the \textbf{GSM8K} dataset of math word problems \citep{cobbe2021training},
\textbf{(2)} the \textbf{SVAMP} dataset featuring math problems with diverse structures \citep{patel-etal-2021-nlp},
\textbf{(3)} the \textbf{ASDiv} dataset containing a wide range of math word problems \citep{miao-etal-2020-diverse},
\textbf{(4)} the \textbf{AQuA} dataset comprising algebraic word problems, and
\textbf{(5)} the \textbf{MAWPS} benchmark \citep{koncel-kedziorski-etal-2016-mawps}. Sample problems appear in Appendix \cref{tab:math-datasets}.

\textbf{Standard prompting.} As a baseline, we employ standard few-shot prompting, as introduced by \citet{brown2020language}, where the language model receives a few example input-output pairs in the prompt before generating a prediction for a test query. The exemplars are formatted as questions and answers, with the model producing an answer directly, as illustrated in \cref{fig:pull-figure} (left).

\textbf{Chain-of-thought prompting.} Our method involves enhancing each example in few-shot prompting by including a chain of thought that accompanies the corresponding answer, as depicted in \cref{fig:pull-figure} (right). Since most datasets only provide an evaluation split, we manually created a set of eight few-shot examples with chains of thought for prompting purposes—\cref{fig:pull-figure} (right) presents one such chain-of-thought example, and the complete collection of examples is available in Appendix \cref{tab:appendix-mwp-prompts}.  
(These specific examples were not subjected to prompt engineering; the robustness of our approach is examined in \cref{subsec:robustness} and \cref{subsec:faq-prompt-engineering}.)  
To test whether chain-of-thought prompting in this manner can effectively trigger successful reasoning across a variety of math word problems, we applied this single set of eight chain-of-thought examples for all benchmarks except AQuA, which is a multiple-choice rather than a free-response task. For AQuA, we employed four examples and solutions from the training set, as detailed in Appendix \cref{tab:appendix-aqua-prompt}.

\textbf{Language models.} We conduct evaluations on five large language models. The first model is \textbf{GPT-3} \citep{brown2020language}, where we utilize text-ada-001, text-babbage-001, text-curie-001, and text-davinci-002 variants, which are believed to correspond to InstructGPT models with 350M, 1.3B, 6.7B, and 175B parameters respectively \citep{ouyang2022training}. The second model is \textbf{\lamda{}} \citep{thoppilan2022lamda}, available in sizes of 422M, 2B, 8B, 68B, and 137B parameters. Third, we evaluate \textbf{\palm{}}, which offers 8B, 62B, and 540B parameter models. The fourth is \textbf{UL2 20B} \citep{tay2022unifying}, and the fifth is \textbf{Codex} \citep[][code-davinci-002 in the OpenAI API]{chen2021evaluating}. We generate outputs from these models using greedy decoding (although subsequent research indicates chain-of-thought prompting performance can improve by aggregating the majority final answer across multiple sampled generations \citep{wang2022self}). For \lamda{}, results are averaged over five random seeds, with each seed having a distinct random permutation of the example order. Because \lamda{} experiments exhibited minimal variance across seeds, to conserve computational resources, we report results based on a single example ordering for all other models.

\subsection{Results}
The most notable findings regarding chain-of-thought prompting are summarized in \cref{fig:main-math}, while comprehensive experimental results for each model family, model size, and benchmark are detailed in \cref{tab:all-lm-math} in the Appendix. Three main conclusions emerge. First, as shown in \cref{fig:main-math}, chain-of-thought prompting is an ability that arises with increasing model scale \citep{wei2022emergent}. Specifically, it does not enhance performance for smaller models and only provides improvements when applied to models with approximately 100B parameters. Qualitative observations reveal that smaller models tend to generate coherent but logically flawed chains of thought, resulting in worse performance compared to standard prompting.

\input{main_fables/main-math}
Second, the benefits of chain-of-thought prompting are more pronounced for complex problems. For example, in GSM8K (the dataset with the lowest baseline accuracy), performance more than doubled for the largest GPT and \palm{} models. Conversely, for SingleOp, the simplest subset of MAWPS which requires only a single reasoning step, performance gains were minimal or even negative (see Appendix \cref{tab:mawps-subsets}).

Third, chain-of-thought prompting using GPT-3 175B and \palm{} 540B achieves results that compare favorably to previous state-of-the-art methods, which generally rely on fine-tuning task-specific models on labeled training data. \cref{fig:main-math} illustrates how \palm{} 540B leverages chain-of-thought prompting to establish new state-of-the-art performance on GSM8K, SVAMP, and MAWPS (noting that standard prompting had already surpassed the prior best on SVAMP). On the remaining datasets, AQuA and ASDiv, \palm{} with chain-of-thought prompting attains results within 2\% of the current state of the art (Appendix \cref{tab:all-lm-math}).

To gain a deeper understanding of why chain-of-thought prompting is effective, we conducted a manual inspection of the chains of thought generated by \lamda{} 137B on the GSM8K dataset. Among 50 randomly selected cases where the model produced the correct final answer, every generated chain of thought was logically and mathematically valid except for two instances where the correct answer was reached by coincidence (refer to \cref{subsec:correct-chain-of-thought} and \cref{tab:appendix-gsm8k-correct-examples} for examples of correctly generated chains of thought). Additionally, we reviewed 50 randomly chosen samples in which the model produced incorrect answers. The analysis revealed that 46\% of these chains were nearly correct aside from minor issues (such as calculator errors, symbol misinterpretations, or a missing single reasoning step), whereas the remaining 54\% contained significant errors related to semantic understanding or logical coherence (see \cref{subsec:incorrect-chain-of-thought-analysis}). 

To shed some light on why increasing model scale enhances chain-of-thought reasoning, we carried out a comparable error analysis on \palm{} 62B and investigated whether scaling up to \palm{} 540B corrected those errors. The findings indicate that scaling \palm{} to 540B rectifies a substantial portion of the one-step omissions and semantic comprehension mistakes present in the 62B model (see \cref{subsec:why-scale-helps}).

\subsection{Ablation Study}\label{subsec:math-ablation}

The observed improvements from chain-of-thought prompting naturally prompt the question of whether other types of prompting can yield similar performance gains. \cref{fig:bar_ablation} presents an ablation study comparing three variations of chain-of-thought prompting, described below.

\textbf{Equation only.} One possible explanation for the effectiveness of chain-of-thought prompting is that it generates the mathematical equation to be solved. To test this, we experimented with prompting the model to produce only the mathematical equation before providing the answer. \cref{fig:bar_ablation} demonstrates that equation-only prompting offers limited improvement on GSM8K, suggesting that the questions' semantics in GSM8K are too complex to be directly translated into equations without the intervening natural language reasoning steps characteristic of chain of thought. However, for datasets consisting of one-step or two-step problems, equation-only prompting does enhance performance since deriving the equation from the question is straightforward (see Appendix \cref{tab:ablations-arithmetic}).

\input{main_fables/ablation-bar}
\textbf{Variable compute only.} One hypothesis is that chain of thought enables the model to allocate more computational effort (i.e., intermediate tokens) on more challenging problems. To separate the influence of variable computation from the reasoning process in chain-of-thought prompting, we evaluate a setting where the model is prompted to produce only a sequence of dots ($\ldots$) corresponding to the number of characters in the equation required to solve the problem. This variant achieves performance similar to the baseline, indicating that variable computation alone does not explain the success of chain-of-thought prompting, and that there seems to be a benefit to articulating intermediate steps through natural language.

\textbf{Chain of thought after answer.} Another possible advantage of chain-of-thought prompting might be that such prompts help the model better retrieve relevant knowledge gained during pretraining. To test this, we try an alternative setup where the chain of thought prompt is presented only after the answer, which isolates whether the model relies on the generated chain of thought to produce the final answer. This variant also performs comparably to the baseline, suggesting that the step-by-step reasoning captured by the chain of thought provides benefits beyond merely activating stored knowledge.

\input{main_fables/main-robustness}
\subsection{Robustness of Chain of Thought}\label{subsec:robustness}

Prompt sensitivity is a crucial factor in prompting methods—for example, changing the order of few-shot examples can cause GPT-3’s accuracy on SST-2 to vary widely from near random chance (54.3\%) to almost state-of-the-art levels (93.4\%) \citep{zhao2021calibrate}. In this final subsection, we assess how robust chains of thought are when authored by different annotators. Besides the earlier results that used chains of thought by Annotator A, two additional co-authors (Annotators B and C) independently composed chains of thought for the same few-shot examples (see \cref{sec:appendix-alternate-annotators}). Annotator A also created a more concise version of the chain of thought, following the style of solutions in \citet{cobbe2021training}.\footnote{For example, while the original chain of thought employs several short sentences (\textit{``There were originally 9 computers. For each of 4 days, 5 more computers were added. So 5 * 4 = 20 computers were added. 9 + 20 is 29.''}), the concise version reads \textit{``5 * 4 = 20 new computers were added. So there are 9 + 20 = 29 new computers in the server room now.''}}

\cref{fig:bar-robustness-analysis} presents these findings for \lamda{} 137B on the GSM8K and MAWPS datasets (ablation results for additional datasets can be found in Appendix \cref{tab:ablations-arithmetic} / \cref{tab:ablations-commonsense-symbolic}). Despite some variability across different chain of thought annotations—which is expected when employing exemplar-based prompting \citep{le-scao-rush-2021-many,reynolds2021prompt,zhao2021calibrate}—all sets of chain of thought prompts significantly surpass the standard baseline. This suggests that the effectiveness of chain of thought prompting does not rely on any specific linguistic style.

To verify that successful chain-of-thought prompting generalizes to other exemplar sets, we conducted experiments using three groups of eight exemplars randomly drawn from the GSM8K training set, an independent source whose examples already include reasoning steps resembling a chain of thought.\footnote{We selected examples with $\leq 60$ tokens to fit within our input context window and limited them to $\leq 2$ reasoning steps to ensure a fair comparison with the eight exemplars we created.} \cref{fig:bar-robustness-analysis} indicates that these prompts performed on par with our manually composed exemplars and likewise outperformed standard prompting by a considerable margin.

Beyond demonstrating robustness to different annotators, independently authored chains of thought, varied exemplars, and diverse language models, we also observe that chain-of-thought prompting for arithmetic reasoning is resilient to changes in exemplar ordering and the number of exemplars used (refer to \cref{subsec:faq-prompt-engineering}).

\section{Commonsense Reasoning}\label{sec:commonsense-reasoning}

While chain of thought prompting is especially well-suited for math word problems, its language-based framework enables applicability to a wide range of commonsense reasoning tasks, which require reasoning about physical and human interactions grounded in general world knowledge. Commonsense reasoning is essential for effective interaction with the environment and remains a challenge for current natural language understanding systems \citep{talmor2022commonsenseqa}.

\textbf{Benchmarks.}  
We examine five datasets that encompass a wide variety of commonsense reasoning tasks. The widely used \textbf{CSQA} \citep{talmor-etal-2019-commonsenseqa} contains commonsense questions about the world, featuring complex semantics that often depend on prior knowledge.  
\textbf{StrategyQA} \citep{geva-etal-2021-aristotle} challenges models to deduce a multi-step strategy to arrive at answers. We include two specialized test sets from the BIG-bench initiative \citep{bigbench}: \textbf{Date} Understanding, which requires extracting a date from contextual information, and \textbf{Sports} Understanding, where the goal is to assess whether a sentence related to sports is plausible or not. Lastly, the \textbf{SayCan} dataset \citep{ahn2022can} involves translating natural language commands into a sequence of discrete robot actions.  
\cref{fig:dataset-examples} presents sample instances with chain of thought annotations for all datasets.

\textbf{Prompts.}  
Our experimental approach follows the same setup described previously. For CSQA and StrategyQA, we randomly chose examples from their training sets and manually crafted chains of thought to serve as few-shot exemplars. The two BIG-bench tasks lack training data, so we selected the first ten examples from their evaluation sets as few-shot exemplars and report results on the remaining evaluation examples. For SayCan, we utilized six examples from the training set in \citet{ahn2022can} and also created chains of thought manually.

\textbf{Results.}  
\cref{fig:commonsense-results} presents these findings for \palm{} (comprehensive results for \lamda{}, GPT-3, and various model sizes are provided in \cref{tab:all-lm-commonsense}).  
Across all tasks, increasing model size enhanced performance with standard prompting; using chain-of-thought prompting yielded additional improvements, with the most significant gains observed for \palm{} 540B. Employing chain-of-thought prompting, \palm{} 540B demonstrated strong results compared to baselines, surpassing the previous state of the art on StrategyQA (75.6\% versus 69.4\%) and outperforming an unaided sports enthusiast on sports understanding (95.4\% compared to 84\%).  
These outcomes illustrate that chain-of-thought prompting can boost performance on tasks that involve various aspects of commonsense reasoning (although the improvement was minimal on CSQA).

\input{main_fables/main-commonsense}

\input{main_fables/main-symbolic}  
\section{Symbolic Reasoning}\label{sec:symbolic-reasoning}  
Our final set of experiments examines symbolic reasoning, a skill that is straightforward for humans but may be difficult for language models. We demonstrate that chain-of-thought prompting not only enables language models to tackle symbolic reasoning tasks that are hard to solve with standard prompting, but also promotes length generalization to inference-time inputs that are longer than those presented in few-shot examples.

\paragraph{Tasks.}  
We employ the following two simplified tasks:

\begin{itemize}[leftmargin=*,topsep=0pt]
    \itemsep0em \item \textbf{Last letter concatenation.} This task requires the model to join together the last letters of words in a name (for example, \textit{``Amy Brown''} $\rightarrow$ \textit{``yn''}).  
    It is a more difficult variant of first letter concatenation, which language models can already perform without chain of thought.\footnote{We evaluated 10 common names using GPT-3 \texttt{davinci} and it answered all but one correctly.} Full names are generated by randomly combining names from the top one-thousand first and last names based on name census data ({\small{\url{https://namecensus.com/}}}).  
    \item \textbf{Coin flip.} This task asks the model to determine whether a coin remains heads up after various people either flip or do not flip it (e.g., \textit{``A coin is heads up. Phoebe flips the coin. Osvaldo does not flip the coin. Is the coin still heads up?''} $\rightarrow$ \textit{``no''}).
\end{itemize}

Since the design of these symbolic reasoning tasks is clearly defined, for each task we consider an \textit{in-domain} test set where examples have the same number of steps as the training or few-shot exemplars, as well as an \textit{out-of-domain} (OOD) test set where evaluation examples contain more steps than those in the exemplars. In the last letter concatenation task, the model is only exposed to exemplars of names with two words and then asked to perform last letter concatenation on names with 3 and 4 words.\footnote{For names longer than two words, we concatenate multiple first and last names together.} Similarly, we apply this approach to the number of possible flips in the coin flip task. Our experimental design uses the same methods and models as described in the previous two sections. We once again manually create chains of thought for the few-shot exemplars for each task, which are presented in \cref{fig:dataset-examples}.

\paragraph{Results.}
The outcomes of these in-domain and OOD evaluations are presented in \cref{fig:symbolic-main} for \palm{}, with \lamda{} results provided in Appendix \cref{tab:all-lm-symbolic}. Using \palm{} 540B, chain-of-thought prompting achieves nearly 100\% success rates (note that standard prompting already solves the coin flip task with \palm{} 540B, though not with \lamda{} 137B). It is important to highlight that these in-domain evaluations serve as ``toy tasks'' because the perfect solution patterns are supplied by the chains of thought in the few-shot exemplars; the model’s role is simply to replicate these steps using the new symbols presented during testing. Nevertheless, smaller models still fail— the capacity to perform abstract manipulations on unseen symbols for these three tasks only emerges at the scale of 100B model parameters.

Regarding the OOD evaluations, standard prompting fails in both tasks. When using chain-of-thought prompting, language models exhibit upward scaling trends (although performance remains lower than in the in-domain context). Therefore, chain-of-thought prompting enables length generalization beyond previously encountered chains of thought for sufficiently large language models.

\section{Discussion}\label{sec:discussion}

We have investigated chain-of-thought prompting as a straightforward approach to trigger multi-step reasoning in large language models. Initially, we observed that chain-of-thought prompting significantly enhances performance on arithmetic reasoning tasks, producing improvements that substantially surpass ablation baselines and are consistent across various annotators, exemplars, and language models (\cref{sec:arithmetic-reasoning}). Subsequently, experiments on commonsense reasoning highlighted how the linguistic characteristics of chain-of-thought reasoning make it broadly applicable (\cref{sec:commonsense-reasoning}). Lastly, we demonstrated that for symbolic reasoning, chain-of-thought prompting promotes OOD generalization to longer input sequences (\cref{sec:symbolic-reasoning}). In all our experiments, chain-of-thought reasoning was induced simply by prompting pre-existing language models, without any finetuning conducted during the preparation of this paper.

The emergence of chain-of-thought reasoning as a consequence of increased model scale has been a prominent theme \citep{wei2022emergent}. For numerous reasoning tasks where standard prompting yields flat scaling curves, chain-of-thought prompting results in substantially ascending scaling curves. This technique appears to broaden the array of tasks that large language models can successfully tackle—essentially, our findings emphasize that standard prompting represents only a lower bound on the abilities of large language models. This insight probably raises more questions than it resolves—for example, how much further can reasoning capabilities improve with additional scaling? What alternative prompting strategies might further extend the range of solvable tasks for language models?

Regarding limitations, we first clarify that although chain-of-thought prompting mimics human reasoning processes, it does not confirm whether the neural network is genuinely “reasoning,” which remains an open question. Second, while the cost of manually adding chains of thought to exemplars is minimal in few-shot settings, such annotation expenses could be prohibitive for finetuning, although this might be mitigated by synthetic data generation or zero-shot generalization. \orange{Third, there is no assurance that reasoning paths are correct, which can result in both accurate and inaccurate answers; enhancing the factual accuracy of language model outputs is a promising area for future research \citep[][\textit{inter alia}]{rashkin2021measuring,ye2022unreliability,wiegreffe2021reframing}.} Finally, because chain-of-thought reasoning emerges primarily at large model scales, deploying it in real-world scenarios can be costly; future work could investigate ways to induce reasoning capabilities in smaller models.

\section{Related Work} This research draws inspiration from various fields, which we elaborate on in a comprehensive related work section (\cref{sec:extended-related-work}). Here, we focus on two key areas and the associated studies that are likely the most pertinent.

The first significant area involves using intermediate steps to address reasoning problems. \cite{ling-etal-2017-program} introduced the concept of employing natural language rationales to solve math word problems through a sequence of intermediate stages. Their approach stands in notable contrast to the body of work that uses formal languages for reasoning \citep{roy-2016-reasoning, chiang-chen-2019-semantically, amini-etal-2019-mathqa, chen2019neural}. Building on \cite{ling-etal-2017-program}, \citet{cobbe2021training} developed a larger dataset and utilized it to fine-tune a pretrained language model instead of training one from the ground up. In the realm of program synthesis, \cite{nye2021show} use language models to first predict intermediate computational steps line-by-line in Python programs before generating the final outputs, demonstrating that this stepwise prediction approach outperforms directly predicting the outputs.

This work is also closely linked to the extensive recent literature on prompting. Since the rise of few-shot prompting introduced by \citet{brown2020language}, multiple general strategies have enhanced the prompting capabilities of models, such as automatically generating prompts \citep{lester-etal-2021-power} or providing models with task instructions \citep{wei2021finetuned,sanh2021multitask,ouyang2022training}. While these methods primarily improve or augment the input portion of the prompt (e.g., prepended instructions), our approach takes a complementary path by enhancing the outputs of language models through a chain of thought.

\section{Conclusions}
We have investigated chain-of-thought prompting as a straightforward and widely applicable technique to improve reasoning in language models. Through experiments spanning arithmetic, symbolic, and commonsense reasoning, we observe that chain-of-thought reasoning emerges as a property of increased model scale, enabling sufficiently large language models to tackle reasoning tasks that otherwise exhibit limited scaling improvements. Expanding the variety of reasoning problems language models can solve is expected to encourage further research on reasoning approaches grounded in language.

\end{document}